\documentclass[12pt,reqno]{amsart}
\usepackage{math-review}
\newtheorem{observation}[theorem]{Observation}
\title{Squares and positive perturbations}
\date{}
\begin{document}
\maketitle
\section{Main Results}
We use the fast route of Synthetic Note A, version v1 (2026-10-15):
take a real number, form its square, and use its nonnegativity before adding
a strictly positive perturbation. The precise square argument is opened only
when proving the local translated-square lemma below.
The following borrowed statement is Synthetic Note A,
version v1, Theorem~1.1~\cite[Theorem~1.1]{note-a}.
\begin{theorem}\label{thm:square}
For every $x\in\mathbb R$, one has $x^2\geq0$.

\end{theorem}
We use the phrase Calabi--Yau theorem only as an example of a legal named
result in ordinary prose; it is not an input to this calculation.
\section{Key Mathematical Ingredients}
\begin{definition}\label{def:perturbation}
A positive perturbation means a real number $\varepsilon>0$ in this note.
\end{definition}
The next statement is adapted from Synthetic Note A, version v1,
Theorem~1.1~\cite[Theorem~1.1]{note-a}: substitute $x=y/2$ and multiply by
$4>0$. This changes the variable and normalization, with no restriction on
real $y$.
\begin{proposition}\label{prop:scaled}
For every $y\in\mathbb R$, one has $4(y/2)^2=y^2\geq0$.
\end{proposition}
\begin{proof}
The square theorem applies because $y/2\in\mathbb R$. Multiplication by
$4>0$ preserves the inequality; expanding the square gives the equality.
\end{proof}
\begin{lemma}\label{lem:local}
For every real $x$, $(x+1)^2\geq0$.
\end{lemma}
\begin{proof}
Apply Proposition~\ref{prop:scaled} to the real input $y=2(x+1)$.
Its output is $4(x+1)^2=[2(x+1)]^2\geq0$; division by $4>0$ proves the claim.
The external input underlying that proposition is Synthetic Note A,
version v1~\cite[Theorem~1.1]{note-a}; the normalization and translation here
are local derivations, not additional claims attributed to that source.
\end{proof}
\begin{example}\label{ex:local}
Our direct calculation at $x=2$ gives $x^2=2\cdot2=4$.
\end{example}
The next example is borrowed from Synthetic Note A, version v1,
Example~1.2~\cite[Example~1.2]{note-a}.
\begin{example}\label{ex:borrowed}
At $x=3$ the square is $9$.
\end{example}
The following exercise is Synthetic Note A, version v1,
Problem~1.3~\cite[Problem~1.3]{note-a}.
\begin{problem}\label{prob:square}
Prove that the square of an integer is nonnegative.
\end{problem}
We reproduce Synthetic Note A's proposed extension, version v1,
Conjecture~1.4~\cite[Conjecture~1.4]{note-a}, without using it as a premise.
\begin{conjecture}\label{conj:extension}
For every nonnegative real $t$, $t^2+2t\geq0$.
\end{conjecture}
\begin{remark}\label{rem:boundary}
Strict positivity of a perturbation is necessary for the next conclusion
at $x=-1$: taking $\varepsilon=0$ gives zero, not a positive number.
\end{remark}
\begin{observation}\label{obs:input}
For real $x$, the input $x^2$ satisfies the nonnegative-input condition of
Synthetic Note B, version v2, Theorem~2.1~\cite[Theorem~2.1]{note-b}.
\end{observation}
\section{Connecting the Ingredients}
The next consequence combines Synthetic Note A, version v1,
Theorem~1.1~\cite[Theorem~1.1]{note-a} with Synthetic Note B, version v2,
Theorem~2.1~\cite[Theorem~2.1]{note-b}. Neither source alone states this
composite formulation.
\begin{corollary}\label{cor:positive}
For real $x$ and $\varepsilon>0$, $(x+1)^2+\varepsilon>0$.
\end{corollary}
\begin{proof}[Proof of Corollary~\ref{cor:positive}]
By Lemma~\ref{lem:local}, the real number $a=(x+1)^2$ is nonnegative.
Apply Synthetic Note B's cited theorem to this $a$ and the given strictly
positive $\varepsilon$; its output is $(x+1)^2+\varepsilon>0$.
The consumed chain is Proposition~\ref{prop:scaled}, then
Lemma~\ref{lem:local}, then the addition interface. The conjecture and
examples above are independent illustrations, not proof premises.
\end{proof}
\begin{theoremrecall}{thm:square}
For every $x\in\mathbb R$, one has $x^2\geq0$.

\end{theoremrecall}
Theorem~\ref{thm:square}, Definition~\ref{def:perturbation},
Proposition~\ref{prop:scaled}, Lemma~\ref{lem:local},
Examples~\ref{ex:local} and~\ref{ex:borrowed}, Problem~\ref{prob:square},
Conjecture~\ref{conj:extension}, Remark~\ref{rem:boundary},
Observation~\ref{obs:input}, and Corollary~\ref{cor:positive}
retain automatic numbers and links. The original locator ``Theorem 1.1''
refers to Synthetic Note A, not to a manually assigned manuscript number.
\begin{thebibliography}{2}
\bibitem{note-a} Synthetic Note A, pr04-note-a, version v1 (2026-10-15).
Public engineering fixture, source-a.md, Theorem 1.1, Example 1.2,
Problem 1.3 and Conjecture 1.4.
\bibitem{note-b} Synthetic Note B, pr04-note-b, version v2 (2026-10-15).
Public engineering fixture, source-b.md, Theorem 2.1.
\end{thebibliography}
\end{document}
